\subsection{Physical-model study: finite-time gains and certified robustness domains}
\label{sec:physicalresults}
\textbf{Status on 12 October 2026.}
A constrained four-state Arrhenius network now meets both predeclared
finite-time benchmarks throughout the simultaneous $\pm5\%$ parameter
box. Continuous interval certificates establish at least $10\%$
KL-distance advantage on $0.4\leq t/\tau\leq0.8$ and at least $5\%$
time savings for every KL threshold in $[10^{-3},10^{-2}]$.
The whole box also lies inside the declared physical parameter ranges;
see Section~\ref{sec:physicalboxcertificate}. This is a separate
physical-model study: two triplets satisfy the original test, so it does
not strengthen the paper's all-triplet-failure counterexample.

\paragraph{Model chosen before screening.}
Let the energies be $E=(0,g_1,g_1+g_2,g_1+g_2+g_3)$, with basins
$A=\{1,4\}$ and $B=\{2,3\}$. Their energy levels are interleaved.
Use a common attempt frequency $\nu_0$ and
\[
 B_{ij}=\max(E_i,E_j)+h_{ij},\qquad
 h_{14}=h_A,\quad h_{23}=h_B,\quad
 h_{12}=h_{13}=h_{24}=h_{34}=h_X.
\]
The three shared activation excesses represent two types of intrabasin
rearrangement and a more costly change of basin. They were not fitted
independently on six edges or reconstructed from a chosen eigenvector.
The dated model specification is \path{research/physical_model_design.txt}.
This is a phenomenological network of metastable configurations; it is
not yet derived from a specified material or continuum potential.
Section~\ref{sec:thermalcollision} supplies an exact stochastic
energy-conserving collision realization of these same rates, with its
bath and contact assumptions stated explicitly.

The refined parameters, rounded here for readability, are
\[
\begin{split}
(g_1,g_2,g_3)&=(0.830894205,\ 0.116352749,\ 2.729136959),\\
(h_A,h_B,h_X)&=(2.114063920,\ 2.822657579,\ 5.640080424),\\
(T_c/T_b,T_h/T_b)&=(3.869407740,\ 9.472085162).
\end{split}
\]
The certified model uses the full decimal inputs in
\path{research/physical_rate_refined.json}, not the rounded display.
Its energy span is $3.67638\,k_BT_b$ and its directed rate ratio is
approximately $584.94$. The smallest and largest directed rates are
approximately $0.000206426\,\nu_0$ and $0.120746\,\nu_0$;
$\tau=147.232\,\nu_0^{-1}$. All rates are strictly positive and detailed
balance follows from the symmetric saddles.

\paragraph{Mechanism and physical interpretation.}
The nominal decay rates in attempt-frequency units are approximately
$0.0067920$, $0.116138$ and $0.130813$. Thus the two faster modes decay
about 17 and 19 times as quickly as the slow mode.
Normalize the slow left eigenfunction to have first component one.
Numerically it is $(1,-1.24304,-1.24253,0.92447)$, nearly constant within
each basin. The cold and hot preparations put respectively $0.46591$
and $0.47963$ of their probability in basin $A$, compared with $0.55459$
at the bath. The hotter preparation therefore has less of the slow
basin-population mismatch, although it starts with a larger total KL
distance. Its additional disturbance within the basins relaxes quickly.
These numbers support the interpretation; exact lumpability is not assumed.
The relevance of metastable populations is established in prior work,
including \href{https://arxiv.org/abs/2304.06420v2}{Biswas, Rajesh and Pal
(2023), Section IV and Fig.~2}. Their piecewise-linear Langevin model
relates initial well populations to relaxation. Their Section VI also
exhibits the effect without a metastable minimum: metastability is not a
general necessary condition. The citation supplies a mechanism precedent,
not a derivation of our four-state rates, an exact lumping theorem, or our
robustness certificate. Neither that general mechanism nor Arrhenius
kinetics is claimed new here.

The barriers are only $2.11$--$5.64\,k_BT_b$ above the higher endpoint.
Consequently a Kramers-type asymptotic rate approximation could have
material errors even with a clear separation of the network's eigenvalues.
The 17-fold spectral separation is between network modes, not a measurement
of intrawell relaxation. A microscopic derivation must establish that
separate scale and quantify the coarse-graining error. No experiment,
absolute clock calibration, preparation cost or continuum validation is
claimed by this calculation.
The separate collision realization defines an exact thermalizing model
without a large-barrier approximation; it does not validate a material's
continuum escape-rate law.

\paragraph{Finite-time evidence.}
Direct 80-digit matrix exponentials give the following times, in
$\nu_0^{-1}$ units:
\begin{center}
\begin{tabular}{rrrrr}
\toprule
KL threshold & Cold time & Hot time & Saved time & Saving\\
\midrule
$10^{-2}$ & $40.2978$ & $27.6018$ & $12.6960$ & $31.5055\%$\\
$10^{-2.5}$ & $125.0952$ & $105.4404$ & $19.6548$ & $15.7118\%$\\
$10^{-3}$ & $209.9825$ & $190.3277$ & $19.6548$ & $9.3602\%$\\
\bottomrule
\end{tabular}
\end{center}
These spot checks are supplemented by outward interval arithmetic.
The certificate proves at least $8.12\%$ time savings for
\emph{every} threshold in $[10^{-3},10^{-2}]$, with both initial
distances above the band and the hot preparation initially farther away.
On the fixed time interval $[58.8927,117.7854]\,\nu_0^{-1}$ (endpoints
rounded here), it proves $D_h\leq0.9D_c$ and $D_c>0.00349$.
An interval LDL inertia calculation brackets the slow eigenvalue and
checks that this interval is longer than $0.25\tau$ and ends before
$5\tau$. The more precise endpoints are in the certificate receipt.
The nominal gain is therefore resolved at finite distances and finite
times, rather than inferred solely from a slow-mode amplitude.

\paragraph{What the interval certificate establishes.}
The checker uses uniformization with $\rho=1/2$:
\[
 e^{Mt}p=e^{-\rho t}\sum_{k=0}^{\infty}
 \frac{(\rho t)^k}{k!}(I+M/\rho)^kp.
\]
The transition matrix is nonnegative and stochastic for the exact input.
Each component of its iterates lies in $[0,1]$, so a geometric upper bound
on the remaining Poisson mass encloses the truncated tail. All exponential
and logarithmic evaluations use 60-digit outward interval arithmetic.
Relative entropy decreases along either relaxation. On a time cell
$[t_i,t_{i+1}]$, it therefore suffices to prove
$D_h(t_i)<0.9D_c(t_{i+1})$. On a threshold cell
$[\delta_i,\delta_{i+1}]$, it suffices to bound
$1-t_h(\delta_i)/t_c(\delta_{i+1})$ from below.
Eighty time cells and sixty threshold cells pass; interval sign checks
bracket each threshold root. This is an executable numerical certificate
with stated mathematical dependencies, not a new Lean formalization.

\paragraph{Robustness: discovery evidence and continuous verification.}
The refined candidate was checked at all $2^8=256$ corners and 3,000
seeded interior points of the independent $\pm5\%$ box in the three
energy gaps, three barrier excesses and two temperature excesses.
At 31 thresholds per point, the smallest observed time saving was
$5.1203\%$. The shortest sampled distance-advantage window exceeded
$1.08\tau$. These samples found no failure of either performance target.
Those samples alone do not exclude an interior parameter or threshold failure.
The subsequent continuous certificates are distinguished from the samples
in Section~\ref{sec:physicalboxcertificate}.
Three further seeded adversarial searches varied all eight uncertain
parameters and the threshold continuously. None found a saving below
$5.1203\%$. These local/global numerical optimization attempts strengthen
the screening evidence but do not prove the global minimum.
The physical range constraints are interval-certified for
the whole box: the energy span is below $3.86021$, every gap exceeds
$0.11053$, barrier excesses lie in $[2.00836,5.92209]$, the rate ratio
is below $993.753$, and both preparations remain in the declared range.

An initial nominal optimization gave about $16.76\%$ minimum sampled time
savings, but 441 of 1,256 subsequent perturbation cases failed at least
one target or range constraint. A second search guarding the physical
box and individual parameter changes still failed in 50 cases.
Three refinements against adverse combinations produced the present
candidate. Retaining these failures is important: checking one parameter
at a time would have overstated robustness. The first two searches used
19,728 and 14,973 objective evaluations; the three refinements used
9,977, 10,094 and 9,743. These bounded searches do not establish global
optimality or a no-go result for other models.

The numerical original-triplet signs are $(F,G)=(-,-),(+,+),(+,+),(-,-)$
for $123,124,134,234$, respectively; the positive margins are large
compared with numerical error. Thus this example is outside the
all-triplet-failure premise. Both finite-time benchmarks have now been certified on the full
$\pm5\%$ box, independently of that additional triplet premise. The physical
approximation check appears in Section~\ref{sec:activatedintermediates};
it does not establish a continuum or material realization.
Scripts and exact input provenance are listed in \path{RUNME.txt}.

The nominal total-variation and Pearson chi-square trajectories are
retained alongside KL in \path{output/figures/physical_finite_time.csv}.
On the sampled grid in $0.4\leq t/\tau\leq0.8$, their distance advantages
are approximately $12.50\%$ and $23.43\%$, respectively, compared with
$23.39\%$ for KL. These diagnostics support the qualitative reversal for
the other distances at the nominal point; the continuous robustness
certificates below concern KL only.

